\section{Discussion: Ambition Without Evidential Inflation}
\subsection{What a Scalar Cannot Say Alone}
A weighted score encodes a preference over the observations and proxies supplied to it. It does not reveal all tradeoffs unless its components are inspectable. Two candidates with the same total may have substantially different profiles. A high value may coexist with an unfavorable comparison in a low-weight component. A score can also improve because a proxy or placeholder changes rather than because physical behavior improves.

The appropriate scientific response is to retain the component vector, source identity and model scope alongside the scalar. This does not prevent continuous ranking. It makes ranking interpretable. The requested ambition of ever-lower loss should therefore be understood as a convention for reporting preference within a stable version, not as evidence that every engineering property improves monotonically or that unlimited compute guarantees an optimum.

\subsection{Claims About Strength, Mass, and Practical Use}
The retained data include modeled mass and mechanical responses. These are valuable observations, but their practical interpretation remains bounded by model assumptions, manufacturing state and missing validation. A low CAD mass does not establish the mass of a finished assembly, and a modeled safety factor does not establish reliability under an unvalidated use scenario. The historical application should not be used to imply operational suitability from a selected numerical example.

This distinction also matters for broader applications. A representation successful in one geometry family may be poorly suited to another. A material label such as titanium does not encode process history, orientation, porosity, surface condition or inspection outcome. Transfer to another domain requires a new statement of what is represented, what is measured and what evidence supports the intended claim.

\subsection{Manufacturing Freedom and Precision}
Metal additive manufacturing can make geometries that are difficult to create through conventional routes. Geometric freedom, however, is different from dimensional precision. A digital model can specify arbitrarily many decimal places without the manufactured object achieving those dimensions. Thermal behavior, surface formation, postprocessing and measurement all influence the relationship between CAD and part.

NIST's geometric test-artifact work provides a concrete reason to measure process capability rather than infer it from the digital model \cite{nistartifact}. The paper therefore treats unprecedented tolerance as a hypothesis requiring metrology, not a conclusion from optimization. A future study would need defined dimensional characteristics, calibrated measurement methods, uncertainty estimates and a comparison baseline. Nominal CAD fidelity alone cannot answer that question.

\subsection{Threats to Validity}
This is a retrospective single-repository study. The retained cohorts were selected by prior workflows and are not random samples. Some source files are unavailable. Historical model and objective revisions limit cross-wave comparison. Shared observations create correlated outcomes, and a large number of labels can exaggerate apparent corroboration. The new publication code verifies projections of retained bytes but cannot establish the truth of every original computation.

The paper also lacks an independent reconstruction experiment performed by a separate research group. Its author red-team review can expose obvious claim and reporting errors, but it does not replace peer review. No controlled agent-assisted representation experiment or physical manufacturing study has been completed for this manuscript. These limitations constrain the current contribution to system explanation, taxonomy and reproducible reporting results.

\subsection{What Would Strengthen a Tier-One Submission}
A stronger empirical submission would add a preregistered research question and an independently reproduced result on a benign reference problem. For the reporting contribution, useful evidence would include reconstruction success, deliberate provenance defects and reviewer comprehension under controlled conditions. For agent assistance, useful evidence would compare typed proposals under fixed independent criteria, retaining invalid and rejected attempts. These are prospective software/research studies, not completed results.

The manuscript should not be submitted as an experimental breakthrough merely because it is long, detailed or visually polished. A high-quality research contribution needs a claim that is both interesting and supported by a method capable of falsifying it. The existing platform's breadth makes such studies plausible; its documentation should make the remaining evidence needs explicit enough that another researcher can carry them out responsibly.

\subsection{Why the Engineering Depth Still Matters}
Careful qualification of claims should not make the platform seem trivial. A linked architecture spanning representation, CAD identity, multiple numerical models, convergence semantics, separate solvers, evidence custody and adversarial reporting controls is substantial engineering work. Making those connections visible lets readers appreciate the system for what it actually contains, including the retained failures and corrections that make its limitations inspectable.

The long-term opportunity is to reduce the cost of asking serious design questions without reducing the quality of the evidence required to answer them. Compute expands exploration. Agents can help manage representational complexity. Verification and measurement constrain interpretation. The value emerges from their coordination, not from substituting any one of them for the others.
